\begin{tomtat}
	\begin{minipage}[t]{0.3\textwidth}
		\vspace{0pt}
		\begin{tomtat}
			\begin{center}
				\textbf{ĐỀ 30}
			\end{center}
			
		\end{tomtat}
	\end{minipage}
	\hfill
	\begin{minipage}[t]{0.5\textwidth}
		\vspace{0pt}
		\begin{tomtat}
			\begin{center}
				\includegraphics[scale=0.08]{images/30}
			\end{center}
		\end{tomtat}
	\end{minipage}
\end{tomtat}
\textbf{I. PHẦN TRẮC NGHIỆM}
\begin{ex}
	Tìm nghiệm của phương trình: $x-7=-2$ là:
	\choice{$-9$}{$5$}{$9$}{$-5$}
	\loigiai{
		Ta có: $x-7 = -2$
		$\Leftrightarrow x = -2+7$
		$\Leftrightarrow x = 5$.
		\\ \textbf{Đáp Án B}
	}
\end{ex}

\begin{ex}
	Tìm nghiệm của bất phương trình: $x-8 \le 0$ là:
	\choice{$x \le 8$}{$x < 8$}{$x \le -8$}{$x \ge 8$}
	\loigiai{
		Ta có: $x-8 \le 0$
		$\Leftrightarrow x \le 8$.
		\\ \textbf{Đáp Án A}
	}
\end{ex}

\begin{ex}
	Tam giác DEF vuông tại D. Có $\widehat{F}=30^\circ$ và cạnh $EF=18$cm. Độ dài cạnh DE là:
	\choice{$9\sqrt{3} \text{ cm}$}{$6\sqrt{3} \text{ cm}$}{$9 \text{ cm}$}{$18\sqrt{3} \text{ cm}$}
	\loigiai{
		Trong tam giác DEF vuông tại D, ta có:
		$\sin F = \frac{DE}{EF}$ (cạnh đối / cạnh huyền)
		$\sin 30^\circ = \frac{DE}{18}$
		$\frac{1}{2} = \frac{DE}{18}$
		$DE = 18 \cdot \frac{1}{2} = 9$ (cm).
		\\ \textbf{Đáp Án C}
	}
\end{ex}

\begin{ex}
	Rút gọn biểu thức: $P=\sqrt{50}+\sqrt{2}$ kết quả là:
	\choice{$5\sqrt{2}$}{$6\sqrt{2}$}{$10$}{Không có đáp án D}
	\loigiai{
		Ta có:
		\begin{align*} P &= \sqrt{50} + \sqrt{2} \\ &= \sqrt{25 \cdot 2} + \sqrt{2} \\ &= 5\sqrt{2} + \sqrt{2} \\ &= (5+1)\sqrt{2} \\ &= 6\sqrt{2} \end{align*}
		\\ \textbf{Đáp Án B}
	}
\end{ex}

\begin{ex}
	Tìm giá trị của a để đồ thị hàm số $y=ax^2$ đi qua điểm $Q(2;4)$.
	\choice{$\frac{1}{2}$}{$\frac{1}{8}$}{$2$}{$1$}
	\loigiai{
		Đồ thị hàm số $y=ax^2$ đi qua điểm $Q(2;4)$, nên tọa độ của Q thỏa mãn phương trình hàm số:
		$4 = a(2)^2$
		$4 = a \cdot 4$
		$a = \frac{4}{4}$
		$a = 1$.
		\\ \textbf{Đáp Án D}
	}
\end{ex}
\begin{ex}
	Công ty Điện lực thống kê lượng điện tiêu thụ (đơn vị: kWh) của một số hộ gia đình ở một khu vực trong tháng. Dữ liệu được ghi lại ở bảng tần số ghép nhóm như sau:
	\begin{center}
		\begin{tabular}{|c|c|c|c|c|}
			\hline
			[100;130) & [130;160) & [160;190) & [190;210) \\
			\hline
			5 & 8 & 7 & 7 \\
			\hline
		\end{tabular}
	\end{center}
	Số hộ gia đình dùng điện từ 130kWh đến dưới 190kWh là:
	\choice{$15$}{$7$}{$20$}{$27$}
	\loigiai{
		Số hộ gia đình dùng điện từ 130kWh đến dưới 190kWh bao gồm các hộ trong nhóm [130;160) và nhóm [160;190).
		Tần số của nhóm [130;160) là 8.
		Tần số của nhóm [160;190) là 7.
		Tổng số hộ gia đình dùng điện từ 130kWh đến dưới 190kWh là: $8 + 7 = 15$.
		\\ \textbf{Đáp Án A}
	}
\end{ex}

\begin{ex}
	Một quả bóng rổ có đường kính 24cm. Tính diện tích bề mặt của quả bóng (lấy $\pi \approx 3,14$; kết quả làm tròn đến hàng đơn vị).
	\choice{$3617 \text{ cm}^2$}{$75 \text{ cm}^2$}{$1809 \text{ cm}^2$}{$904 \text{ cm}^2$}
	\loigiai{
		Đường kính của quả bóng rổ là $d = 24$ cm.
		Bán kính của quả bóng rổ là $R = \frac{d}{2} = \frac{24}{2} = 12$ cm.
		Diện tích bề mặt của hình cầu được tính bằng công thức $S = 4\pi R^2$.
		$S = 4\pi (12)^2 = 4\pi \cdot 144 = 576\pi \text{ (cm}^2\text{)}$.
		Với $\pi \approx 3,14$:
		$S \approx 576 \cdot 3,14 = 1808,64 \text{ (cm}^2\text{)}$.
		Làm tròn kết quả đến hàng đơn vị, ta được $S \approx 1809 \text{ cm}^2$.
		\\ \textbf{Đáp Án C}
	}
\end{ex}

\begin{ex}
	Cho 2 điểm A và B thuộc (O; 10cm) sao cho $\widehat{AOB} = 54^\circ$ (lấy $\pi \approx 3,14$). Độ dài cung nhỏ AB (làm tròn đến hàng đơn vị) là:
	\choice{$18 \text{ cm}$}{$9 \text{ cm}$}{$5 \text{ cm}$}{$10 \text{ cm}$}
	\loigiai{
		Độ dài cung tròn được tính bằng công thức $l = \frac{\pi R n}{180}$, trong đó R là bán kính và n là số đo góc ở tâm (tính bằng độ).
		Với $R=10$ cm và $n=54^\circ$, ta có:
		$l = \frac{\pi \cdot 10 \cdot 54}{180} = \frac{540\pi}{180} = 3\pi \text{ (cm)}$.
		Với $\pi \approx 3,14$:
		$l \approx 3 \cdot 3,14 = 9,42 \text{ (cm)}$.
		Làm tròn kết quả đến hàng đơn vị, ta được $l \approx 9$ cm.
		\\ \textbf{Đáp Án B}
	}
\end{ex}

\textbf{II. PHẦN TỰ LUẬN}
\begin{ex}
	Cho $a>0$ và $a \neq 25$. Rút gọn biểu thức $Q = \left( \frac{1}{\sqrt{a}+5} + \frac{1}{\sqrt{a}-5} \right) : \left( 1 + \frac{5}{\sqrt{a}-5} \right)$.
	\loigiai{
		Điều kiện: $a>0, a \neq 25$.
		Ta có:
		\begin{align*} Q &= \left( \frac{(\sqrt{a}-5) + (\sqrt{a}+5)}{(\sqrt{a}+5)(\sqrt{a}-5)} \right) : \left( \frac{\sqrt{a}-5+5}{\sqrt{a}-5} \right) \\ &= \frac{2\sqrt{a}}{a-25} : \frac{\sqrt{a}}{\sqrt{a}-5} \\ &= \frac{2\sqrt{a}}{(\sqrt{a}-5)(\sqrt{a}+5)} \cdot \frac{\sqrt{a}-5}{\sqrt{a}} \end{align*}
		Vì $a>0$ và $a \neq 25$, nên $\sqrt{a} \neq 0$ và $\sqrt{a}-5 \neq 0$.
		\begin{align*} Q &= \frac{2\cancel{\sqrt{a}}}{\cancel{(\sqrt{a}-5)}(\sqrt{a}+5)} \cdot \frac{\cancel{(\sqrt{a}-5)}}{\cancel{\sqrt{a}}} \\ &= \frac{2}{\sqrt{a}+5} \end{align*}
		Vậy $Q = \frac{2}{\sqrt{a}+5}$.
	}
\end{ex}

\begin{ex}
	Biết rằng phương trình bậc hai $x^2+4x-3=0$ có hai nghiệm phân biệt $x_1, x_2$. Không giải phương trình, tính giá trị của biểu thức $R = -x_1-x_2+x_1^2x_2+x_1x_2^2$.
	\loigiai{
		Phương trình $x^2+4x-3=0$ có $a=1, b=4, c=-3$.
		$\Delta' = b'^2 - ac = (2)^2 - 1(-3) = 4+3=7 > 0$.
		Phương trình có hai nghiệm phân biệt $x_1, x_2$.
		Theo hệ thức Vi-ét, ta có:
		$$ \begin{cases} x_1 + x_2 = -\frac{b}{a} = -\frac{4}{1} = -4 \\ x_1 x_2 = \frac{c}{a} = \frac{-3}{1} = -3 \end{cases} $$
		Ta biến đổi biểu thức R:
		\begin{align*} R &= -(x_1+x_2) + x_1x_2(x_1+x_2) \end{align*}
		Thay giá trị từ hệ thức Vi-ét vào biểu thức R:
		\begin{align*} R &= -(-4) + (-3)(-4) \\ &= 4 + 12 \\ &= 16 \end{align*}
		Vậy $R=16$.
	}
\end{ex}

\begin{ex}
	Bạn Nam gieo con xúc xắc một lần. Tính xác suất của biến cố E: "Số chấm xuất hiện trên mặt con xúc xắc là số nguyên tố".
	\loigiai{
		Khi gieo con xúc xắc một lần, không gian mẫu $\Omega = \{1; 2; 3; 4; 5; 6\}$.\\
		Số phần tử của không gian mẫu là $n(\Omega) = 6$.\\
		Biến cố E: "Số chấm xuất hiện trên mặt con xúc xắc là số nguyên tố".\\
		Các số nguyên tố trong các số từ 1 đến 6 là: $\{2; 3; 5\}$.\\
		Vậy, $E = \{2; 3; 5\}$.\\
		Số phần tử của biến cố E là $n(E) = 3$.\\
		Xác suất của biến cố E là:\\
		$P(E) = \frac{n(E)}{n(\Omega)} = \frac{3}{6} = \frac{1}{2}$.
	}
\end{ex}

\begin{ex}
	Giải hệ phương trình: $\begin{cases} -2x+y=4 \\ x+y=1 \end{cases}$.
	\loigiai{
		Ta có hệ phương trình:
		$$ \begin{cases} -2x+y=4 \quad &(1) \\ x+y=1 \quad &(2) \end{cases} $$
		Trừ vế theo vế của phương trình (1) cho phương trình (2), ta được:
		\begin{align*} (-2x+y) - (x+y) &= 4-1 \\ -2x+y-x-y &= 3 \\ -3x &= 3 \\ x &= -1 \end{align*}
		Thay $x=-1$ vào phương trình (2):
		\begin{align*} -1+y &= 1 \\ y &= 1+1 \\ y &= 2 \end{align*}
		Vậy nghiệm của hệ phương trình là $(x;y) = (-1;2)$.
	}
\end{ex}
\begin{ex}
	Bà Bình gửi 30\,000\,000 đồng vào một ngân hàng theo kì hạn 1 năm. Sau 1 năm bà Bình không rút tiền lãi về mà số tiền lãi đó được cộng vào tiền vốn để tính lãi cho kì hạn tiếp theo. Tổng số tiền bà Bình nhận được sau 2 năm là 33\,708\,000 đồng. Hỏi lãi suất ngân hàng mà bà Bình gửi là bao nhiêu? Biết trong 2 năm đó lãi suất ngân hàng không thay đổi.
	\loigiai{
		Gọi lãi suất hàng năm của ngân hàng là $x$ (dưới dạng số thập phân).
		Điều kiện: $x > 0$.
		
		Số tiền vốn ban đầu của bà Bình là $P_0 = 30\,000\,000$ đồng.
		
		Sau 1 năm, số tiền lãi bà Bình nhận được là: $L_1 = P_0 \cdot x = 30\,000\,000x$ (đồng).
		Tổng số tiền bà Bình có sau 1 năm (cả vốn lẫn lãi) là:
		$P_1 = P_0 + L_1 = P_0 + P_0x = P_0(1+x) = 30\,000\,000(1+x)$ (đồng).
		
		Số tiền này được dùng làm vốn cho năm thứ hai.
		Số tiền lãi bà Bình nhận được trong năm thứ hai là: $L_2 = P_1 \cdot x = 30\,000\,000(1+x)x$ (đồng).
		Tổng số tiền bà Bình nhận được sau 2 năm là:
		$P_2 = P_1 + L_2 = P_1 + P_1x = P_1(1+x) = P_0(1+x)(1+x) = P_0(1+x)^2$.
		
		Theo đề bài, tổng số tiền bà Bình nhận được sau 2 năm là 33\,708\,000 đồng, nên ta có phương trình:
		\begin{align*} 30\,000\,000(1+x)^2 &= 33\,708\,000 \\ (1+x)^2 &= \frac{33\,708\,000}{30\,000\,000} \\ (1+x)^2 &= \frac{33708}{30000} \\ (1+x)^2 &= \frac{11236}{10000} \end{align*}
		Ta thấy $11236 = 106^2$ và $10000 = 100^2$.
		\begin{align*} (1+x)^2 &= \left(\frac{106}{100}\right)^2 \\ (1+x)^2 &= (1,06)^2 \end{align*}
		Vì $1+x > 0$ (do $x>0$), ta có:
		\begin{align*} 1+x &= 1,06 \\ x &= 1,06 - 1 \\ x &= 0,06 \end{align*}
		Lãi suất ngân hàng là $0,06 \times 100\% = 6\%$.
		
		Vậy, lãi suất ngân hàng mà bà Bình gửi là 6\% một năm.
	}
\end{ex}
\begin{ex}
	Cho tam giác ABC nhọn, nội tiếp đường tròn (O). Các đường cao AD, BE, CF cắt nhau tại H.
	\begin{enumerate}
		\item[a)] Chứng minh tứ giác BFEC nội tiếp đường tròn.
		\item[b)] Kẻ đường kính AK. Gọi S là điểm đối xứng của K qua C, L là điểm đối xứng với K qua B. Chứng minh: $AF \cdot AB = AE \cdot AC$ và $DL=DS$.
	\end{enumerate}
\end{ex}

\begin{ex}
	\begin{enumerate}
		\item[a)] Một cửa sổ dạng hình vòm gồm phần hình chữ nhật phía dưới và nửa hình tròn phía trên. Phần hình chữ nhật có chiều dài của cạnh đứng là 1m chiều dài cạnh ngang là 1,2m. Biết giá mỗi mét vuông cửa là 700\,000 đồng. Hãy tính giá tiền của cửa sổ trên?
		\item[b)] Cho 3 số dương a, b, c thỏa mãn $abc = 1$. Chứng minh:
		$$ \frac{1}{a^3+b^3+1} + \frac{1}{b^3+c^3+1} + \frac{1}{c^3+a^3+1} \le 1 $$
	\end{enumerate}
\end{ex}

\setcounterref{ex}{0}
